\chapter{Reproducing the Results}\label{app:repro}

The repository's \texttt{README.md} and \texttt{docs/how-it-works.md} describe setup in detail. On
Linux, \texttt{./setup.sh} creates a Python environment with ns-3 and SUMO. The main commands are:

\begin{center}\small
\begin{tabular}{@{}lp{9.5cm}@{}}
\toprule
Command & Result \\
\midrule
\texttt{make test} & all automated tests \\
\texttt{make live} & one ns-3 run narrated in the terminal \\
\texttt{make viz} & the same run as a live map in the browser \\
\texttt{make scenario-3obu}, \texttt{-low}, \texttt{-veryhigh} & demonstration scenarios: one normal vehicle, one attacker and one emergency vehicle; 30 vehicles; 300 vehicles \\
\texttt{make evaluate-all} & the sweeps behind Chapter~\ref{ch:results} (about a day on 8 cores) \\
\texttt{make emergency} & the emergency-vehicle sweep \\
\texttt{make detector-benchmark} & the VeReMi Extension benchmark \\
\texttt{make ledger-bench} & the Fabric Raft and SmartBFT benchmark \\
\texttt{make fabric-up}, \texttt{make gateway} & a real Fabric network and its REST gateway \\
\bottomrule
\end{tabular}
\end{center}

Every table, figure and number in Chapters~\ref{ch:results} and in the abstract is generated from
the experiment outputs by \texttt{python experiments/thesis\_data.py}, which writes the folder
\texttt{docs/report/generated}. Each run folder contains \texttt{config.json} with every option of the
run, so any result can be repeated exactly. Every run also has its \texttt{metrics.csv} with one row per
event, from which all metrics are computed.
